\begin{emphbox}
    \begin{equation*}
        (\Omega, \mathbb F, P)
    \end{equation*}
\end{emphbox}

\subsection{Grundlagen}

\begin{sectionbox}
    \subsubsection{Ergebnis}
    Ausgang, welchen ich messe/vorfinde/beobachte nach Durchführung des Experiments.

    \subsubsection{Ergebnismenge $\Omega$}
    Alle möglichen Ausgänge eines Zufallsexperiments.
\end{sectionbox}

\begin{sectionbox}
    \subsubsection{Ereignismenge ($\sigma$-Algebra)}
    Teilmengen $A$ von $\Omega$. Ausgänge, die uns interessieren (z.B. Augenzahl gerade) \\

    \textbf{Partition} Disjunkte Teilmengen, welche zusammen $\Omega$ ergeben (TODO: Mitschrift)

    \textbf{Komplement}
    $A^C = \Omega\setminus A \qquad A \cup A^C = \Omega$ \\

    \textbf{Überlappen von verschiedenen Ereignismengen} \\
    Falls $A\cap B = \emptyset$ hat man 8 verschieden Fälle. Falls $A\cap B \ne \emptyset$ hat man 16 (TODO: Mitschrift) \\

    $\mathbb F$ ist eine konsistente Ereignismenge, wenn gilt: (TODO)
    \begin{itemize}
        \item $\Omega \in \mathbb F$
        \item $A \in \mathbb F \Rightarrow A^C \in \mathbb F$
        \item $A_i\in\mathbb F \Rightarrow A_1\cup A_2 \cup ... \in \mathbb F$
    \end{itemize}

\end{sectionbox}

\begin{sectionbox}
    \subsubsection{De-Morgan}
    $(A^C \cup B^C)^C=A\cap B$
\end{sectionbox}


\begin{sectionbox}
    \subsubsection{Wahrscheinlichkeitsmaß $P$}
    $P$ ist eine Funktion $P : \mathbb F \rightarrow [0,1] \subset \mathbb R$. \\

\end{sectionbox}